Effective communication depends on a fixed cadence and on a single place where decisions are recorded. Every recurring activity below produces an artefact -- a meeting note, a reviewed pull request, an issue, or an updated document -- so that a decision made in a conversation does not have to be remembered by the people who attended it.

\begin{table}[H]
\centering
\small
\setlength{\tabcolsep}{4pt}
\renewcommand{\arraystretch}{1.3}
\begin{tabular}{|p{3.1cm}|p{2.7cm}|p{3.9cm}|p{2.1cm}|p{2.6cm}|}
\hline
\rowcolor{fptHeader}
\textbf{Communication Item} &
\textbf{Who/ Target} &
\textbf{Purpose} &
\textbf{When, Frequency} &
\textbf{Type, Tool, Method(s)} \\
\hline

Daily stand-up &
Development team &
Report progress since the previous stand-up, state the next item, and raise blockers &
Daily, 15 minutes &
Discord voice channel, short meeting note \\
\hline

Backlog refinement &
Development team &
Re-estimate candidate items, split oversized items, and confirm dependencies &
Weekly &
GitHub project board, Discord discussion \\
\hline

Sprint Planning &
Team members, supervisor &
Confirm the sprint goal and the commitment up to the effort cap &
Start of each sprint &
Google Meet, GitHub issues \\
\hline

Sprint Review &
Team members, supervisor, invited users &
Demonstrate the increment on a running environment and collect feedback &
End of each sprint &
Google Meet, live demonstration on the staging environment \\
\hline

Sprint Retrospective &
Development team &
Agree at most two process changes and verify the previous sprint's changes &
End of each sprint &
Discord, retrospective note in the documentation folder \\
\hline

Design and schema review &
Team members, module owner &
Agree on the data model, the API contract, and the permission entry before implementation &
Before each high-complexity item &
Google Meet, design note committed to the repository \\
\hline

Pull request review &
Author and at least one reviewer &
Check behaviour, error handling, authorisation, and test coverage &
Per change &
GitHub pull request review \\
\hline

Defect triage &
Team members &
Classify defects by severity and assign an owner &
Twice per sprint &
GitHub issues with severity labels \\
\hline

Weekly progress report &
Team leader, supervisor &
Report status against the sprint cap, the risk register, and any support required &
Weekly &
Written report sent by email, Google Meet when a discussion is needed \\
\hline

\end{tabular}
\caption{Project Communication Plan}
\label{tab:communication-plan}
\end{table}

Communication is conducted through Google Meet for synchronous reviews and supervisor meetings, Discord for day-to-day coordination and quick technical discussion, GitHub for code, pull request discussion, and the issue list, and Google Drive for documents, reports, and shared figures. Figma is used for interface design hand-off, where comments on a frame are resolved before the corresponding frontend item is started.

Meeting outcomes are recorded in a short note with the date, the participants, the decisions taken, and the action items with an owner. The note is stored in the shared documentation folder and linked from the task board, so a member who missed the meeting can recover the decision without asking. A decision that changes scope, schedule, or an existing interface is also recorded in the relevant document rather than only in the meeting note, and the change is announced in the next stand-up.
